\chapter{Nhật ký làm việc nhóm}

\section{Kế hoạch làm việc nhóm}

Quá trình thực hiện đồ án được chia thành bốn giai đoạn chính. Giai đoạn đầu tập trung khởi tạo dự án Maven, thống nhất cấu trúc package, bộ khung unit test và giao diện cơ sở. Giai đoạn hai, mỗi thành viên triển khai độc lập các use case được phân công trên nhánh (branch) riêng. Giai đoạn ba tiến hành tích hợp toàn bộ luồng nghiệp vụ end-to-end, từ bước tạo yêu cầu nhập hàng cho đến khi đối soát hàng về kho. Giai đoạn cuối cùng dành cho việc sửa lỗi (bug fix), hoàn thiện bộ kiểm thử, viết tài liệu báo cáo và đóng gói sản phẩm cuối cùng.

\begin{longtblr}{
  colspec = {Q[0.16\textwidth, l] Q[0.32\textwidth, l] X[l]},
  hlines,
  vlines,
  rowhead = 1,
}
\textbf{Giai đoạn} & \textbf{Mục tiêu} & \textbf{Kết quả cần có} \\
Phase 1 & Khởi tạo nền kiến trúc & Maven project, JDK 22, cấu trúc package, test skeleton, README build/run. \\
Phase 2 & Triển khai từng use case & Code và unit test cho UC001--UC007 trên nhánh cá nhân. \\
Phase 3 & Tích hợp end-to-end & Luồng tạo yêu cầu, hỏi tồn kho, phân bổ, đặt hàng, nhận hàng chạy liền mạch. \\
Phase 4 & Ổn định và nộp bài & Build, test, package, tài liệu báo cáo và ảnh minh họa hoàn chỉnh. \\
\end{longtblr}

\section{Nhật ký họp nhóm}

\begin{longtblr}{
  colspec = {Q[0.16\textwidth, l] Q[0.28\textwidth, l] X[l]},
  hlines,
  vlines,
  rowhead = 1,
}
\textbf{Thời điểm} & \textbf{Nội dung} & \textbf{Kết luận} \\
Tuần 1 & Chốt phạm vi final project & Chọn JavaFX, Maven, JDK 22, Supabase REST; giữ 7 UC theo phân công. \\
Tuần 2 & Rà soát thiết kế và dữ liệu dùng chung & Thống nhất các field lõi: \fieldname{requestId}, \fieldname{siteCode}, \fieldname{merchandiseCode}, \fieldname{quantity}, \fieldname{unit}, \fieldname{desiredDeliveryDate}. \\
Tuần 3 & Tích hợp service và giao diện & Tách màn hình theo use case; dùng \ucname{ImportOrderingFacade} để giảm phụ thuộc trực tiếp. \\
Tuần 4 & Kiểm thử và hoàn thiện báo cáo & Chạy Maven test, kiểm tra luồng demo, hoàn thiện nội dung thiết kế và nhật ký nhóm. \\
\end{longtblr}

\section{Khó khăn và cách giải quyết}

\begin{longtblr}{
  colspec = {Q[0.30\textwidth, l] Q[0.30\textwidth, l] X[l]},
  hlines,
  vlines,
  rowhead = 1,
}
\textbf{Khó khăn} & \textbf{Ảnh hưởng} & \textbf{Cách giải quyết} \\
Dữ liệu giữa các UC phụ thuộc nhau & UC004, UC005, UC006 không chạy nếu thiếu dữ liệu từ UC001, UC002, UC003. & Thiết kế facade và store chung, tích hợp theo thứ tự phụ thuộc. \\
Rule phân bổ có nhiều tiêu chí & Dễ tạo kết quả khác nhau giữa tài liệu và code. & Tách \ucname{AllocationPolicy}, viết test cho ưu tiên tồn kho, ngày nhận và phương thức vận chuyển. \\
Kết nối Supabase cần cấu hình môi trường & Ứng dụng có thể lỗi khi chạy lần đầu. & Đưa biến môi trường vào hướng dẫn chạy và dùng \ucname{InMemoryStore} cho test. \\
Giao diện nhiều màn hình & Dễ trùng logic xử lý form và thông báo lỗi. & Tạo component hỗ trợ như \ucname{FormFieldFactory}, \ucname{StatusBadge}, \ucname{Toast}, \ucname{ConfirmDialog}. \\
Quản lý role và dữ liệu master dễ lấn sang nghiệp vụ chính & UC007 có thể làm rối luồng nhập hàng tuyến tính. & Giữ UC007 trong màn hình quản trị riêng, chỉ cung cấp quyền, user, log và cấu hình. \\
\end{longtblr}

\section{Kết quả đạt được}

Nhóm đã xây dựng thành công phần mềm cốt lõi cho Hệ thống Đặt hàng Nhập khẩu, đáp ứng trọn vẹn các yêu cầu nghiệp vụ đã đề ra. Về mặt kỹ thuật, dự án ứng dụng kiến trúc phân tầng (Layered/DDD) chặt chẽ với giao diện JavaFX, xử lý luồng thao tác qua Application Service, gom nhóm nghiệp vụ tại Domain Model và kết nối cơ sở dữ liệu thật qua Supabase. Về mặt tài liệu, nhóm đã biên soạn hoàn chỉnh và thống nhất các tài liệu cấu thành đồ án bao gồm: Đặc tả yêu cầu (SRS), Phân tích \& Thiết kế Use Case, Thiết kế Giao diện, Thiết kế Chi tiết, cùng với việc giải trình rõ ràng các Nguyên lý và Mẫu thiết kế đã được hiện thực hóa trong mã nguồn.
